
In the case now under consideration, \(X_0,X_1,X_2\) are affine. We
may therefore suppose that
\[
  X_i=\operatorname{Spec}A_i,\qquad
  d_j=\operatorname{Spec}\delta_j,\qquad
  d'_k=\operatorname{Spec}\delta'_k,
\]
where the \(A_i\) are commutative rings and the
\(\delta_j,\delta'_k\) are ring homomorphisms. The diagram
\hyperref[sga3:v:lemma:1-2:equation:0-1-2]{(0,1,2)} may then be replaced by
the following dual diagram:

\begin{equation*}
  \resizebox{0.57\textwidth}{!}{%
    \SGAThreeNativeInput{figures/exp5/p262_dual_cocartesian_012star.png}}
  \phantomsection\label{fig:sga3-exp5-dual-012star}
\end{equation*}
The two left-hand squares in this diagram are cocartesian.

Let \(B\) be the subring of \(A_0\) consisting of the elements
\(a\in A_0\) for which \(\delta_0(a)=\delta_1(a)\).

\medskip
\phantomsection\label{proof:sga3-V-4.1-a}
\noindent\textbf{a)} We first show that \(A_0\) is integral over
\(B\). For \(a\in A_0\), let
\begin{equation}
  P_{\delta_1}(T,\delta_0(a))
  =T^n-\sigma_1T^{n-1}+\cdots+(-1)^n\sigma_n
  \label{eq:sga3-V-4.1-characteristic-polynomial}
\end{equation}
be the characteristic polynomial of \(\delta_0(a)\), where \(A_1\)
is regarded as an \(A_0\)-algebra through \(\delta_1\); see
Bourbaki, \emph{Alg\`ebre}, VIII, \S12, and
\emph{Alg\`ebre commutative}, II, \S5, Exercise~9.
Since the two left-hand squares of the diagram
\hyperref[fig:sga3-exp5-dual-012star]{\((0,1,2)^\ast\)} are
cocartesian, one has
\begin{align}
 \delta_0\!\left(P_{\delta_1}(T,\delta_0(a))\right)
   &=P_{\delta'_2}(T,\delta'_0\delta_0(a)), \\
 \delta_1\!\left(P_{\delta_1}(T,\delta_0(a))\right)
   &=P_{\delta'_2}(T,\delta'_1\delta_0(a)).
\end{align}
Because \(\delta'_0\delta_0=\delta'_1\delta_0\), it follows that
\[
 \delta_0\!\left(P_{\delta_1}(T,\delta_0(a))\right)
 =
 \delta_1\!\left(P_{\delta_1}(T,\delta_0(a))\right);
\]
equivalently,
\[
  \delta_0(\sigma_i)=\delta_1(\sigma_i)
  \qquad\text{for every }i.
\]

On the other hand, the Cayley--Hamilton theorem gives
\[
 \delta_0(a)^n-\delta_1(\sigma_1)\delta_0(a)^{n-1}
 +\cdots+(-1)^n\delta_1(\sigma_n)=0.
\]
Since \(\delta_1(\sigma_i)=\delta_0(\sigma_i)\), one also has
\[
 \delta_0(a)^n-\delta_0(\sigma_1)\delta_0(a)^{n-1}
 +\cdots+(-1)^n\delta_0(\sigma_n)=0.
\]
Hence
\[
  a^n-\sigma_1a^{n-1}+\cdots+(-1)^n\sigma_n=0,
\]
because there is a homomorphism
\(\tau:A_1\to A_0\) such that
\(\tau\delta_0=\operatorname{id}_{A_0}\), so \(\delta_0\) is
injective. It follows that \(A_0\) is integral over \(B\).
